% Modello "Beamer" — LaTeX Editor Offline
% Contenuto originale, liberamente modificabile.
% Richiede la classe "beamer" nella distribuzione TeX installata.
\documentclass[11pt]{beamer}
\usepackage[utf8]{inputenc}
\usepackage[italian]{babel}

\title{Titolo della presentazione}
\author{Nome Autore}
\date{\today}

\begin{document}

\begin{frame}
  \titlepage
\end{frame}

\begin{frame}{Sommario}
  \tableofcontents
\end{frame}

\section{Contesto}

\begin{frame}{Contesto}
  \begin{itemize}
    \item primo punto della presentazione;
    \item secondo punto della presentazione.
  \end{itemize}
\end{frame}

\section{Risultati}

\begin{frame}{Un risultato}
  \begin{equation}
    \label{eq:beamer}
    a^{2} + b^{2} = c^{2}
  \end{equation}
  L'equazione~\eqref{eq:beamer} è un esempio dentro un frame.
\end{frame}

\begin{frame}{Conclusioni}
  \begin{itemize}
    \item sintesi dei risultati;
    \item sviluppi futuri.
  \end{itemize}
\end{frame}

\end{document}
